% 2010 年辽宁卷理科第 18 题（2）①：两张待完成的频率分布直方图坐标系。
% 来源 PDF：普通高考/2010/2010辽宁理.pdf，SHA-256 ba46f70f9d718652c24d14b28a0b473b768e864248ff0a9f1feb2374aec708f9。
% 原页：第 1 页（A3 三栏）第 3 栏下部，书页码 611。
% 重绘依据：照原页画空坐标系（原页即要求学生补画小矩形）——纵轴「频率/组距」自
%   0.01 到 0.08，横轴 60 到 85（每 5 一格），纵横虚线网格、两轴各有一处断裂记号；
%   左为图 1：注射药物 A，右为图 2：注射药物 B。不代画柱形，保持与原页一致。
\begin{tikzpicture}[x=0.7cm, y=0.49cm, semithick]
  % ---------- 图 1：注射药物 A ----------
  \begin{scope}
    % 坐标轴
    \draw[->, >=stealth] (0,0) -- (0,9.15);
    \draw[->, >=stealth] (0,0) -- (7.05,0);
    % 虚线网格：竖线 x=1..6（即 60..85），横线 y=1..8（即 0.01..0.08）
    \foreach \i in {1,2,3,4,5,6} {
      \draw[dashed, dash pattern=on 3pt off 3pt] (\i,0) -- (\i,8);
    }
    \foreach \j in {1,2,3,4,5,6,7,8} {
      \draw[dashed, dash pattern=on 3pt off 3pt] (0,\j) -- (6,\j);
    }
    % 纵轴刻度值
    \foreach \j/\v in {1/0.01, 2/0.02, 3/0.03, 4/0.04, 5/0.05, 6/0.06, 7/0.07, 8/0.08} {
      \node[anchor=east, inner sep=1pt] at (-0.12,\j) {\v};
    }
    % 横轴刻度值
    \node[anchor=north, inner sep=2pt] at (0,-0.12) {0};
    \foreach \i/\v in {1/60, 2/65, 3/70, 4/75, 5/80, 6/85} {
      \node[anchor=north, inner sep=2pt] at (\i,-0.12) {\v};
    }
    % 轴断裂记号
    \draw (0.42,0.10) -- (0.56,-0.10) -- (0.70,0.10);
    \draw (0.62,0.10) -- (0.76,-0.10) -- (0.90,0.10);
    \draw (0.10,0.42) -- (-0.10,0.56) -- (0.10,0.70);
    % 轴名
    \node[align=center, inner sep=1pt] at (0.92,10.75) {频率};
    \draw[semithick] (0.36,10.25) -- (1.48,10.25);
    \node[align=center, inner sep=1pt] at (0.92,9.75) {组距};
    \node[anchor=west, inner sep=2pt] at (7.15,1.0) {疱疹面积};
    % 图题
    \node[anchor=north, inner sep=3pt] at (3,-1.05) {图 1：注射药物 $A$};
  \end{scope}

  % ---------- 图 2：注射药物 B ----------
  \begin{scope}[xshift=7.95cm]
    \draw[->, >=stealth] (0,0) -- (0,9.15);
    \draw[->, >=stealth] (0,0) -- (7.05,0);
    \foreach \i in {1,2,3,4,5,6} {
      \draw[dashed, dash pattern=on 3pt off 3pt] (\i,0) -- (\i,8);
    }
    \foreach \j in {1,2,3,4,5,6,7,8} {
      \draw[dashed, dash pattern=on 3pt off 3pt] (0,\j) -- (6,\j);
    }
    \foreach \j/\v in {1/0.01, 2/0.02, 3/0.03, 4/0.04, 5/0.05, 6/0.06, 7/0.07, 8/0.08} {
      \node[anchor=east, inner sep=1pt] at (-0.12,\j) {\v};
    }
    \node[anchor=north, inner sep=2pt] at (0,-0.12) {0};
    \foreach \i/\v in {1/60, 2/65, 3/70, 4/75, 5/80, 6/85} {
      \node[anchor=north, inner sep=2pt] at (\i,-0.12) {\v};
    }
    \draw (0.42,0.10) -- (0.56,-0.10) -- (0.70,0.10);
    \draw (0.62,0.10) -- (0.76,-0.10) -- (0.90,0.10);
    \draw (0.10,0.42) -- (-0.10,0.56) -- (0.10,0.70);
    \node[align=center, inner sep=1pt] at (0.92,10.75) {频率};
    \draw[semithick] (0.36,10.25) -- (1.48,10.25);
    \node[align=center, inner sep=1pt] at (0.92,9.75) {组距};
    \node[anchor=west, inner sep=2pt] at (7.15,1.0) {疱疹面积};
    \node[anchor=north, inner sep=3pt] at (3,-1.05) {图 2：注射药物 $B$};
  \end{scope}
\end{tikzpicture}
